% =============================================================================
%  RESUMO
% =============================================================================
\begin{resumo}

Este livro apresenta uma jornada autodidata e sistemática através dos fundamentos 
matemáticos, estatísticos e computacionais que sustentam a engenharia de 
ecossistemas cognitivos artificiais, tendo como objeto central de estudo e 
experimentação o \opencodeeco{} (versão 5.4.0, ciclo evolutivo R23). 
A obra estrutura-se em oito capítulos que conduzem o leitor do nível zero 
(fundamentos básicos de matemática e lógica) ao nível PhD (validação 
científica, produção acadêmica \qualisA{} e defesa perante banca examinadora).

O Capítulo 1 estabelece a base matemática e estatística --- álgebra linear, 
cálculo, probabilidade, inferência estatística e teoria da informação --- 
essencial para a compreensão dos algoritmos de inteligência artificial. 
O Capítulo 2 introduz os conceitos fundamentais de IA, aprendizado de máquina, 
redes neurais profundas (transformers) e arquiteturas de agentes autônomos. 
O Capítulo 3 mergulha na arquitetura do \opencodeeco, detalhando o padrão 
três camadas (MCP $\rightarrow$ Skill $\rightarrow$ Agent), a metodologia 
SDD+TDD, o barramento de eventos e o sistema de injeção de dependência.

O Capítulo 4 explora o \textit{Pipeline de Scanners Epistemológicos} --- 
Noológico, Teleológico, Evolutivo, Refinamento e MCSP --- e o sistema de 
metacognição (SPEC-036) com Self-Model N0-N3 e monitor metacognitivo. 
O Capítulo 5 apresenta o Trust Engine (SPEC-038) com Behavioral Gate, 
Natural Forgetting (modelo Atkinson-Shiffrin) e Governança Cooperativa 
(princípios de Ostrom). O Capítulo 6 detalha a Token Economy (SPEC-022/023/024) 
com \textit{staking}, \textit{slashing}, mercado de taxas e trilha de auditoria.

O Capítulo 7 descreve o sistema de experimentação e validação científica: 
benchmark CORA-Eval (150 tarefas, 10 dimensões), validação Aletheia 
(Lean 4), pipeline de produção acadêmica MASWOS com 49 agentes especializados 
e o sistema de correção iterativa Qualis A1. O Capítulo 8, finalmente, 
apresenta o roteiro completo para produção e defesa de dissertação: 
metodologia PPGTE/UFC, simulação de banca com agentes, protocolo de 
anonimato e métricas de avaliação DAP.

Cada capítulo contém definições formais, ilustrações (TikZ), exemplos 
práticos executáveis no \opencodeeco{}, exercícios progressivos (nível 0 
ao PhD), notas de rodapé com citações verificáveis e links ativos para 
fontes primárias. Todas as referências seguem a norma NBR 6023/2023 da 
ABNT, e as citações seguem a NBR 10520/2023.

\textbf{Palavras-chave}: Engenharia de Software. Inteligência Artificial. 
Sistemas Multiagentes. Metacognição. Trust Engine. OpenCode Ecosystem. 
SDD+TDD. Qualis A1. ABNT.

\end{resumo}
